Durch die ehrenwörtliche Erklärung versichern Sie, die gängigen Regeln der
wissenschaftlichen Redlichkeit eingehalten zu haben. Für die Formatierung gilt
im Prinzip dasselbe wie im vorangegangen Abschnitt. Es kommt lediglich ein
Feld für Ihre Unterschrift hinzu. Für die Erklärung wurde dieser Vorlage
deshalb eine separate Umgebung spendiert. Eine ehrenwörtliche Erklärung fügen
Sie wie in Listing~\ref{lst:final:erklaerung} dargestellt ein.

\lstinputlisting[language={[LaTeX]{TeX}},style=block,escapechar=\!,float,caption={Ehrenwörtliche Erklärung},label={lst:final:erklaerung}]{code/final_statement.de.tex}

Das Beispiel entspricht der Formulierung, wie sie zum Zeitpunkt der Verfassung
dieser Anleitung\footnote{April 2017} vom Zentralen Prüfungsamt der TU~Chemnitz
festgelegt wurde. Beachten Sie, dass eine Versicherung an Eides statt
normalerweise nicht notwendig ist.

Sofern die Überschrift der Erklärung anzupassen ist, kann dies über das
optionale Argument der Umgebung geschehen. Ein Unterschriftenfeld erzeugen Sie
durch die Kommandosequenz \verb+\tucsignature+. Sie müssen lediglich Ihren
Namen angeben. Bei Gruppenarbeiten mit mehreren Autoren, können Sie den Befehl
mehrfach nacheinander verwenden. Sollten Sie eine breitere Unterschriftenzeile
benötigen, weil Sie z.\,B. einen längeren Namen haben, geben Sie im optionalen
Argument ein entsprechendes Längenmaß an. Standardmäßig werden \unit[5]{cm}
angesetzt.

Die Umgebung beginnt automatisch eine neue rechte Seite und stellt auch sicher,
dass keine Seitenzahl gedruckt wird.
